\section{Scree Plot Analysis}
\label{appendix-scree-plot}

\begin{figure}[htbp]
  \centering
  \includegraphics[width=\textwidth]{scree_plot_parallel_analysis.png}
  \caption{Scree Plot with Parallel Analysis}
  \label{fig-appendix-scree}
\end{figure}

\clearpage

\section{Detailed Exploratory Factor Analysis}
\label{appendix-fa-detailed}

\begin{figure}[htbp]
  \centering
  \includegraphics[width=\textwidth]{factor_loadings_comparison.png}
  \caption{Factor Loadings Comparison: Closed vs Open Responses Formats}
  \label{fig-appendix-fa}
\end{figure}

\clearpage

\section{Survey Questions Wording}
\label{appendix-wording}

\subsection{Demographic Questions}
\begin{itemize}
  \item How old are you? (Please answer using numbers only)
  \item To which gender do you identify?
  \item In which province or territory do you live?
  \item What is your type of dwelling?
  \item What language do you speak most often at home?
  \item What is the highest level of education that you have completed?
  \item What is the average income of your household before taxes?
\end{itemize}

\subsection{Environment Questions}
\begin{itemize}
  \item How much should the federal government spend on the environment?
  \item Environmental regulation should be stricter, even if it leads to consumers having to pay higher prices.
  \item When there is a conflict between protecting the environment and creating jobs, jobs should come first.
  \item What do you think is the main cause of climate change?
  \item To what extent should the government support the oil industry?
  \item The threat of climate change is exaggerated
  \item The federal carbon tax should be maintained as an environmental measure
\end{itemize}

\subsection{Immigration Questions}
\begin{itemize}
  \item Too many recent immigrants just don't want to fit in Canadian society.
  \item Immigrants take jobs away from Canadians.
  \item How much should the federal government spend on immigrants and minorities?
  \item How many immigrants should the country admit?
  \item Canada's border controls should be strengthened to limit the number of people applying for refugee status
\end{itemize}

\subsection{Vote Choice Question}
\begin{itemize}
  \item Which party did you vote for during the 2025 federal election?
\end{itemize}

\clearpage

\section{Prompt wordings}
\label{appendix-prompts}

\subsection{Master Prompt Wording}
\label{appendix-master-prompt}

{\small
\begin{verbatim}
You are an expert at creating optimized prompts for AI systems that process 
survey data. Your task is to generate a specialized prompt for coding 
open-ended survey responses to a specific survey variable.

## Variable Information:
- Variable name: {variable_name}
- Question text: {question_text}
- Variable type: {variable_type}
- Domain/topic: {variable_domain}
- Sample response values: {sample_values}

## Response Categories/Options:
{categories_info}

## Sample Open-ended Responses (if available):
{sample_responses}

## Language Information:
{language_info}
IMPORTANT: This is a bilingual survey (French and English). Responses may be 
in either language.

## Your Task:
Generate an optimal prompt consisting of TWO parts:

### PART 1: SYSTEM MESSAGE
Create a system message that:
- Defines the AI assistant's role and expertise for this specific variable type
- Explains the task clearly (mapping open responses to codes)
- Provides domain-specific guidance relevant to this variable's topic
- Emphasizes returning only the numeric code
- Includes any special considerations for this variable type
- CRITICAL: Explicitly mentions handling both French and English responses
- Provides key French translations for common responses (oui=yes, non=no, etc.)
- Warns against coding valid French responses as Don't know

### PART 2: USER TEMPLATE
Create a user message template that:
- Uses placeholder variables: {{variable_name}}, {{question_text}}, 
  {{options_block}}, {{open_response}}
- Is formatted clearly for easy reading
- Includes appropriate context for this variable type
- Follows this general structure but adapt the labels/sections as needed:

Variable: {{variable_name}}
Question: {{question_text}}

[Appropriate section title for the options]:
{{options_block}}

Open-ended response:
"{{open_response}}"

## Requirements:
- Be specific to this variable's domain and characteristics
- Consider the types of responses likely for this question
- Optimize for accuracy in mapping responses to the correct codes
- Keep instructions clear and concise
- Ensure the prompt will work well for the sample responses shown
- MUST handle bilingual responses correctly (French and English)
- Include guidance on common French political terms if relevant 
  (e.g., Libéral = Liberal Party)
- Provide clear French-English mappings for agreement/disagreement terms

## Output Format:
IMPORTANT: Return ONLY valid JSON in this exact format. Do not include any 
other text or explanation:

{
  "system_message": "Your system message here...",
  "user_template": "Your user template here..."
}

Generate the optimized prompt now:
\end{verbatim}
}

\clearpage

\subsection{Variable Prompt Wording}
\label{appendix-variable-prompt}

{\small
\begin{verbatim}
### System Message

You are an expert AI assistant specializing in coding survey data for 
political science and public opinion research. Your task is to analyze an 
open-ended survey response and assign the single most appropriate numeric 
code based on the provided categories. The response is for a question about 
federal government spending on the environment.

## Task Instructions:

1. **Analyze the Response**: Carefully read the open-ended response to 
   understand the respondent's opinion on the level of environmental 
   spending.

2. **Identify Sentiment**: Determine if the respondent wants to spend MORE, 
   LESS, or maintain the SAME level of spending.

3. **Assign Code**: Select the single numeric code that best corresponds to 
   this sentiment.

4. **Output**: Return ONLY the single numeric code. Do not provide any 
   explanation, text, or punctuation.

## Critical Considerations for this Variable:
- **Bilingual Survey**: This is a bilingual survey. Responses may be in 
  English or French (FR-CA). You MUST correctly interpret both languages. 
  Do NOT code a valid French response as 'Don't know' (6).
- **Focus on Quantity**: The core of the question is about the *amount* of 
  spending. Look for words indicating an increase, decrease, or maintenance 
  of the current level.
- **Typo in Options**: Note that the option 'Spend about the same as snow' 
  contains a typo and should be interpreted as 'Spend about the same as now'.

## French Language Guidance:
To ensure accuracy, use the following French-to-English translations for 
common terms related to this question:

**More / Increase:**
- `Plus` = More
- `Augmenter` / `Augmentation` = Increase
- `Davantage` = More
- `Beaucoup plus` = A lot more
- `Plus qu'en ce moment` = More than right now

**Less / Decrease:**
- `Moins` = Less
- `Diminuer` / `Réduire` = Decrease / Reduce
- `Beaucoup moins` = A lot less

**Same / Maintain:**
- `Pareil` / `La même chose` = The same
- `Comme maintenant` = Like now
- `Le même montant` = The same amount
- `Garder le même` = Keep the same

**Don't Know / Refusal:**
- `Je ne sais pas` / `Sais pas` / `NSP` = I don't know
- `Aucune idée` = No idea

If the response is ambiguous, expresses no opinion, or is a clear 'don't know' 
or refusal to answer (in either language), use code 6.

### User Template

Variable: {variable_name}
Question: {question_text}

Response Options & Codes:
{options_block}

Open-ended response:
"{open_response}"

### Response Options

1: Spend less
4: Spend about the same as snow
5: Spend more
6: Don't know/Prefer not to answer
\end{verbatim}
}

\clearpage

\subsection{Textual Literacy Prompt Wording}
\label{appendix-literacy-prompt}

{\small
\begin{verbatim}
You are a textual literacy assessor. Analyze the following concatenated survey 
responses from a single participant and provide a literacy score between 0 and 
10, where:

- 0 = literacy level of a severely handicapped child (minimal vocabulary, 
      basic errors, incoherent)
- 10 = literacy level of a Harvard literature professor (sophisticated 
       vocabulary, complex syntax, nuanced expression)

Consider these factors:
- Vocabulary sophistication and range
- Grammatical accuracy and complexity
- Coherence and logical flow
- Spelling and punctuation
- Ability to express nuanced ideas
- Use of appropriate register for the context

Participant's responses (separated by |): "[concatenated_text]"

Return your assessment in this exact JSON format:
{
  "literacy_score": [number between 0 and 10],
  "justification": "[brief explanation of the score]"
}

Important: The score must be a number between 0 and 10 (can include decimals 
like 7.5).
\end{verbatim}
}

\clearpage

\subsection{Strategic Vote Prompt Wording}
\label{appendix-strategic-prompt}

{\small
\begin{verbatim}
Analyze this Canadian federal voting response and extract two pieces of 
information:

Which federal party did the respondent vote for? Use ONLY these party codes: 
lpc, cpc, bq, ndp, gpc, ppc, or other

Does the respondent seem to have voted strategically? (1 for strategic, 0 for 
not strategic)

Response to analyze: response_text

Return your answer in this exact JSON format: 
{ party: [party code], strategic: [0 or 1] }

Important guidelines:

Use ONLY these exact party codes (lowercase):

- lpc = Liberal Party of Canada / Parti libéral du Canada / Libéral / 
        Liberal / PLC
- cpc = Conservative Party of Canada / Parti conservateur du Canada / 
        Conservateur / Conservative / PCC
- bq = Bloc Québécois / Bloc québécois / Bloc / BQ
- ndp = New Democratic Party / Nouveau Parti démocratique / NPD / NDP
- gpc = Green Party of Canada / Parti vert du Canada / Vert / Green / PVC
- ppc = People's Party of Canada / Parti populaire du Canada / PPC
- other = any other party or unclear

Responses may be in French or English

Look for words like strategic, to block, to prevent, par stratégie, pour 
bloquer, pour empêcher for strategic voting

If unclear about strategic voting, default to 0 (not strategic)

If party is unclear, use 'other'
\end{verbatim}
}

\clearpage

\section{Missing Data Analysis}
\label{appendix-missing-data}

\subsection{Missing Data Rates by Item and Response Format}

\begin{table}[htbp]
  \centering
  \begin{tabular}{llcccp{1.5cm}}
    \toprule
    \textbf{Domain} & \textbf{Item} & \textbf{Closed \%} & \textbf{Open \%} & \textbf{$\Delta$} & \textbf{\textit{p}} \\
    \midrule
    Immigration & Gov Spending & 7.8 & 37.8 & +30.0 & <.001 \\
    Immigration & Admit Less & 5.4 & 32.7 & +27.3 & <.001 \\
    Environment & Gov Spending & 3.9 & 28.4 & +24.5 & <.001 \\
    Environment & Anti Oil Support & 2.0 & 9.8 & +7.8 & <.001 \\
    Immigration & Don't Fit In & 4.8 & 10.9 & +6.1 & <.001 \\
    Environment & Priority & 2.1 & 4.7 & +2.6 & <.001 \\
    Environment & Climate Real & 1.5 & 3.9 & +2.4 & <.001 \\
    Environment & Regulation & 1.5 & 3.8 & +2.3 & <.001 \\
    Immigration & Stricter Border & 8.1 & 10.3 & +2.2 & .055 \\
    Immigration & Take Jobs & 1.3 & 3.0 & +1.7 & .003 \\
    Environment & Climate Human & 0.9 & 1.3 & +0.3 & .492 \\
    Environment & Carbon Tax & 10.5 & 10.3 & $-$0.2 & .885 \\
    \bottomrule
  \end{tabular}
  \caption{Missing data rates by survey item and response format}
  \label{table-missing-data}
\end{table}

\noindent\textit{Note:} $\Delta$ = difference in percentage points (Open -- Closed). The three items with highest non-response (>24 percentage points) all request quantitative policy estimates. All tests are chi-square tests of independence.

\subsection{Survey Fatigue Analysis}

\begin{figure}[htbp]
  \centering
  \includegraphics[width=\textwidth]{fig_na_fatigue_analysis.png}
  \caption{Missing data rates by question position and response format}
  \label{fig-appendix-survey-fatigue}
\end{figure}

\noindent\textit{Note:} Missing data rates by question position (1--12) and response format. Lines show linear trends with 95\% confidence bands. Neither format shows significant correlation between position and non-response ($\rho_{\text{closed}}$ = 0.49, \textit{p} = .11; $\rho_{\text{open}}$ = 0.39, \textit{p} = .21), and the negative interaction term ($\beta$ = $-$0.031, \textit{p} = .018) indicates slightly less fatigue progression in open-ended conditions.

\clearpage

\section{Table of Descriptive Statistics}
\label{appendix-descriptive-table}

% latex table generated in R 4.5.2 by xtable 1.8-4 package
% Wed Nov 12 10:08:15 2025
\begin{table}[H]
\centering
\caption{Descriptive Statistics Comparison: Open-ended vs. Closed-ended Treatment. Domain: Env = Environment, Immi = Immigration. Stat: SD = Standard Deviation. Each variable shows three statistics with corresponding Open and Closed treatment values, their ratio (Open/Closed), and difference (Open - Closed).} 
\label{tab:descriptive_statistics}
\begingroup\small
\begin{tabular}{lllrrrr}
  \toprule
Domain & Variable & Stat & Open & Closed & Ratio & Diff \\ 
  \midrule
Env & Against oil support & Mean & 0.773 & 0.769 & 1.01 & 0.004 \\ 
  Env & Against oil support & Median & 0.750 & 1.000 & 0.75 & -0.250 \\ 
  Env & Against oil support & SD & 0.287 & 0.292 & 0.98 & -0.005 \\ 
  Env & Carbon tax (maintained) & Mean & 0.741 & 0.698 & 1.06 & 0.043 \\ 
  Env & Carbon tax (maintained) & Median & 1.000 & 0.667 & 1.50 & 0.333 \\ 
  Env & Carbon tax (maintained) & SD & 0.385 & 0.332 & 1.16 & 0.053 \\ 
  Env & Climate change (human cause) & Mean & 0.899 & 0.900 & 1.00 & -0.001 \\ 
  Env & Climate change (human cause) & Median & 1.000 & 1.000 & 1.00 & 0.000 \\ 
  Env & Climate change (human cause) & SD & 0.301 & 0.300 & 1.00 & 0.001 \\ 
  Env & Climate change (not exaggerated) & Mean & 0.934 & 0.869 & 1.07 & 0.065 \\ 
  Env & Climate change (not exaggerated) & Median & 1.000 & 1.000 & 1.00 & 0.000 \\ 
  Env & Climate change (not exaggerated) & SD & 0.194 & 0.208 & 0.93 & -0.014 \\ 
  Env & Environment priority over jobs & Mean & 0.711 & 0.581 & 1.22 & 0.130 \\ 
  Env & Environment priority over jobs & Median & 0.750 & 0.750 & 1.00 & 0.000 \\ 
  Env & Environment priority over jobs & SD & 0.313 & 0.268 & 1.17 & 0.045 \\ 
  Env & Government spending (more) & Mean & 0.829 & 0.786 & 1.06 & 0.043 \\ 
  Env & Government spending (more) & Median & 1.000 & 1.000 & 1.00 & 0.000 \\ 
  Env & Government spending (more) & SD & 0.329 & 0.325 & 1.01 & 0.004 \\ 
  Env & Regulation (stricter) & Mean & 0.776 & 0.697 & 1.11 & 0.079 \\ 
  Env & Regulation (stricter) & Median & 1.000 & 0.750 & 1.33 & 0.250 \\ 
  Env & Regulation (stricter) & SD & 0.380 & 0.292 & 1.30 & 0.088 \\ 
  Immi & Admit fewer immigrants & Mean & 0.535 & 0.645 & 0.83 & -0.110 \\ 
  Immi & Admit fewer immigrants & Median & 0.500 & 0.750 & 0.67 & -0.250 \\ 
  Immi & Admit fewer immigrants & SD & 0.289 & 0.268 & 1.08 & 0.021 \\ 
  Immi & Immigrants don't fit in & Mean & 0.330 & 0.294 & 1.12 & 0.036 \\ 
  Immi & Immigrants don't fit in & Median & 0.250 & 0.250 & 1.00 & 0.000 \\ 
  Immi & Immigrants don't fit in & SD & 0.314 & 0.281 & 1.12 & 0.033 \\ 
  Immi & Immigrants take jobs & Mean & 0.090 & 0.167 & 0.54 & -0.077 \\ 
  Immi & Immigrants take jobs & Median & 0.000 & 0.000 &  & 0.000 \\ 
  Immi & Immigrants take jobs & SD & 0.208 & 0.208 & 1.00 & 0.000 \\ 
  Immi & Immigration spending (less) & Mean & 0.402 & 0.472 & 0.85 & -0.069 \\ 
  Immi & Immigration spending (less) & Median & 0.500 & 0.500 & 1.00 & 0.000 \\ 
  Immi & Immigration spending (less) & SD & 0.413 & 0.362 & 1.14 & 0.052 \\ 
  Immi & Stricter border control & Mean & 0.493 & 0.592 & 0.83 & -0.099 \\ 
  Immi & Stricter border control & Median & 0.667 & 0.667 & 1.00 & 0.000 \\ 
  Immi & Stricter border control & SD & 0.270 & 0.316 & 0.85 & -0.047 \\ 
   \bottomrule
\end{tabular}
\endgroup
\end{table}

\clearpage

\section{Response Format Effects on Political Attitudes}
\label{appendix-descriptive-effects}

\begin{figure}[htbp]
  \centering
  \includegraphics[width=\textwidth]{response_format_comparison_3x4.png}
  \caption{Density Distribution Comparing Open-Ended vs. Close-Ended Survey Responses}
  \label{fig-appendix-descriptive-effects}
\end{figure}

\clearpage

\section{Subgroups Analysis}
\label{appendix-subgroups}

\begin{table}[htbp]
  \centering
  \begin{tabular}{lcc}
    \toprule
    \textbf{SES Variable} & \textbf{Environment Scale} & \textbf{Immigration Scale} \\
    \midrule
    Age (continuous) & F=1.1, p=0.295 & F=5.87, p=0.016 \\
    Gender (female) & F=1.3, p=0.254 & F=0.42, p=0.515 \\
    Education & F=1.69, p=0.12 & F=1.05, p=0.394 \\
    Income & F=0.59, p=0.762 & F=1.07, p=0.379 \\
    Language & F=0.16, p=0.854 & F=1.28, p=0.278 \\
    Province & F=1.15, p=0.318 & F=1.68, p=0.079 \\
    \bottomrule
  \end{tabular}
  \caption{Robustness checks: Treatment × Subgroup interactions}
  \label{table-robustness}
\end{table}

\noindent\textit{Note:} Non-significant interactions indicate treatment effects are stable across subgroups.

